\chapter*{Introduction}
\label{sec:introduction}
\addcontentsline{toc}{chapter}{Introduction}

Over the past decade, deep learning has evolved from a promising research direction into one of the most transformative technologies in modern science and engineering. Built upon artificial neural networks trained on large volumes of data, deep learning models have achieved, and in several cases surpassed, human-level performance on tasks once considered intractable for machines. Convolutional architectures revolutionized computer vision, enabling accurate image classification, segmentation, and object detection~\cite{oshea2015}; sequence and attention-based models reshaped natural language processing, culminating in large language models capable of fluent reasoning, translation, and code generation~\cite{Goodfellow2016}. Beyond these flagship applications, deep learning now underpins systems for speech recognition, recommendation, medical diagnosis, and autonomous decision-making, establishing itself as a general-purpose framework for learning complex, nonlinear relationships directly from data~\cite{Ruder2016}.

This success has rapidly propagated into the natural sciences, where deep learning has accelerated discovery across disciplines. Neural networks have driven breakthroughs in protein structure prediction, the discovery of new stable materials~\cite{merchant2023}, and medium-range weather forecasting~\cite{lam2023}, and they have become powerful tools in fundamental physics, from the search for new particles to the analysis of gravitational-wave signals~\cite{karagiorgi_2022}. A particularly fertile development has been the emergence of physics-informed learning, in which known physical laws, expressed as conservation principles, symmetries, or differential equations, are embedded directly into the training process. Physics-Informed Neural Networks (PINNs)~\cite{Raissi2019, karniadakis2021} exemplify this paradigm, restricting the space of admissible solutions to those consistent with the governing physics and thereby achieving accurate, data-efficient models even in regimes where purely data-driven approaches would falter. It is this physics-informed perspective that motivates and unifies the work developed in this thesis.

Within physics, quantum information science is an especially natural and demanding domain for these methods. Quantum systems are described by objects, state vectors, density operators, and unitary evolutions, whose dimension grows exponentially with the number of constituents, rendering exact classical treatments intractable beyond modest system sizes~\cite{lloyd1996}. This exponential barrier, together with the abundance of structured data produced by quantum experiments, makes quantum information a setting where the representational power and scalability of deep learning can have substantial impact. Indeed, machine learning has already been applied across a broad spectrum of quantum-information tasks: the reconstruction of quantum states from measurement data, or quantum state tomography~\cite{carrasquilla_qst_2019, danaci2021}; the characterization of quantum processes through quantum process tomography~\cite{huang2025}; the solution of the quantum many-body problem via neural quantum states~\cite{carleo2017, carleo2019, huang2022provably, requena2023}; the learning of $N$-representability conditions in quantum chemistry~\cite{sagersmith2022, delgadogranados2025, shao2023}; and the control of quantum systems through physics-informed networks that solve the underlying differential equations~\cite{norambuena2024}. Across these applications, a consistent theme emerges: deep learning does not merely accelerate existing computations, but learns the latent, low-dimensional structure that often underlies high-dimensional quantum problems.

This thesis contributes to this growing body of work by developing two deep-learning frameworks for the study of physical systems in quantum information, each addressing a problem whose exact treatment is computationally demanding and each guided by the physics-informed principle of embedding known structure into the learning process.

The first problem concerns the dynamics of quantum systems governed by time-dependent Hamiltonians. Unlike the time-independent case, where the time-evolution operator reduces to a simple matrix exponential $e^{-iHt}$, time-dependent Hamiltonians give rise to a time-ordered exponential that precludes a closed-form solution and introduces substantial mathematical and numerical complexity. Series representations such as the Dyson or Magnus expansions involve nested time integrals whose number grows with the expansion order, and the associated numerical errors accumulate over long evolution times~\cite{blanes2009, wiebe_simulating_2011}; even on quantum hardware, simulating such dynamics is generally more demanding than the time-independent case, particularly for rapidly varying or highly oscillatory Hamiltonians~\cite{an2022}. Despite these difficulties, time-dependent Hamiltonians are central to quantum technologies, including quantum control and gate optimization~\cite{Ansel_2024} and the quantum simulation of complex many-body systems~\cite{georgescu2014, cao2025}. A variety of approaches have been developed to address this challenge, including truncated Dyson series~\cite{berry_2015, krieferova_2019}, Magnus-expansion-based algorithms for quantum control~\cite{Dalgaard_2022, chen_quantum_2023, fang_time-dependent_2024}, modified Suzuki--Trotter decompositions for unitary interpolation~\cite{schilling2024}, and tensor-network methods~\cite{vidal_efficient_2004, tepaske2023, causer2024}, as well as machine-learning approaches that learn quantum dynamics or reconstruct time-dependent Hamiltonians~\cite{che_2021, han_2021, franceschetto2025}. In Chapter~\ref{ch:results-pinn-unitaries}, we introduce a physics-informed deep learning framework that, given the time-evolution operator on a finite set of sampled times, interpolates it over the entire time domain, without ever being supplied the Hamiltonian explicitly, by incorporating unitarity as a physical constraint and, for larger systems, predicting an effective Hamiltonian propagated through structure-preserving integrators. This work has been submitted for publication~\cite{guerra2026interpolation}.

The second problem concerns the Quantum Marginal Problem (QMP): given a set of reduced density matrices describing subsystems of a multipartite system, determining whether they are consistent with some global quantum state, and, when they are, reconstructing such a state. This is a fundamental question in quantum theory whose consistency version is complete for the Quantum Merlin--Arthur (QMA) complexity class~\cite{liu2006consistency, broadbent2022qma}, reflecting its computational intractability in the general case. The QMP carries profound implications: it would enable more efficient computation of ground states of local Hamiltonians~\cite{liu2007nrep} and lies at the heart of quantum chemistry, where it appears as the $N$-representability problem~\cite{judd2001coulson}, itself QMA-complete. Special cases have been addressed theoretically, from Klyachko's conditions for the compatibility of one-body marginals~\cite{klyachko2004} to studies on the uniqueness of pure global states determined by their marginals~\cite{linden2002, diosi2004, wyderka2017} and semidefinite-programming formulations of symmetric instances~\cite{aloy2020}. In Chapter~\ref{ch:results-qmp}, we introduce a deep-learning approach that combines a convolutional denoising autoencoder with the Marginal Imposition Operator~\cite{uzcategui2022mio} to reconstruct physically valid global density matrices compatible with a prescribed set of marginals, for systems of up to eight qubits, and demonstrate that, once trained, it offers a faster alternative to state-of-the-art semidefinite programming solvers without compromising solution quality. This work has been published in \textit{Machine Learning: Science and Technology}~\cite{uzcategui2025marginals}.

The remainder of this thesis is organized as follows. Chapters~\ref{sec:linear_algebra}, \ref{sec:quantum_information}, and \ref{sec:deep_learning} establish the necessary background in linear algebra, quantum information, and deep learning, respectively, providing the unified mathematical language in which the two applications are formulated. Chapter~\ref{ch:results-pinn-unitaries} presents the interpolation of unitary time-evolution operators via physics-informed neural networks, and Chapter~\ref{ch:results-qmp} presents the reconstruction of density matrices from quantum marginals. Finally, Chapter~\ref{sec:conclusions} draws together the conclusions of both works, discusses their implications for large-scale deep learning in quantum information, and outlines directions for future research.
